\documentclass[UTF8]{ctexart}
\usepackage{amsmath, amssymb, geometry}
\usepackage[table]{xcolor}
\usepackage{enumitem}
\usepackage{tcolorbox}
\usepackage{tikz}
\usetikzlibrary{calc, arrows.meta}
\usepackage{graphicx}
\geometry{a4paper, left=2cm, right=2cm, top=2.5cm, bottom=2.5cm}

% 定义颜色
\definecolor{bookblue}{RGB}{0, 155, 230}
\definecolor{grassgreen}{RGB}{196, 214, 160}
\definecolor{weekendred}{RGB}{200, 50, 50}
\definecolor{solargreen}{RGB}{50, 150, 50}

% 定义红色环境
\newenvironment{redenv}{\color{red}}{}

% 定义详细解析环境
\newenvironment{detailedanalysis}{%
    \color{red}
    \begin{enumerate}[label=\textbf{第\arabic*步：}, leftmargin=*, align=left, itemsep=1.5ex]
}{%
    \end{enumerate}
}

% 定义概念框
\newenvironment{conceptbox}{%
    \begin{tcolorbox}[colback=red!5, colframe=red!50!black, title=概念解析, fonttitle=\bfseries]
}{%
    \end{tcolorbox}
}

% 定义公式推导环境
\newenvironment{formuladerivation}{%
    \begin{enumerate}[label={\textbf{公式\arabic*：}}, leftmargin=*, align=left]
}{%
    \end{enumerate}
}

% 定义题目和解析的分隔线
\newcommand{\problemrule}{\noindent\rule{\textwidth}{1.5pt}\vspace{-0.8em}}
\newcommand{\analysisrule}{\noindent\rule{\textwidth}{0.8pt}\vspace{-0.8em}}

% 去掉默认段首缩进
\setlength{\parindent}{0pt}

\begin{document}

% 标题
\begin{center}
    {\LARGE\bfseries\color{bookblue} 习题 2.5 详细讲义}\\[0.5em]
    {\large 一元二次方程应用题专题解析}
\end{center}

\vspace{1em}

% ========== 第1题开始 ==========
\problemrule
\section*{【题目1】年平均增长率问题}
某市政府为落实保障性住房政策，去年已投入 3 亿元资金，并规划投入资金逐年增加，明年将投入 12 亿元资金用于保障性住房建设. 求这两年中投入资金的年平均增长率.

\begin{redenv}
\analysisrule
\subsection*{逐步解析}

\begin{detailedanalysis}
\item \textbf{问题本质分析}

这是一道典型的\textbf{增长率问题}。题目涉及三个时间点：
\begin{itemize}
    \item \textbf{去年}：初始投入 $a = 3$ 亿元（基准量）
    \item \textbf{今年}：经过一年增长后的投入
    \item \textbf{明年}：经过两年增长后的投入 $b = 12$ 亿元（终值）
\end{itemize}

\begin{conceptbox}
\textbf{平均增长率的核心概念}：

设初始量为 $a$，年平均增长率为 $r$，经过 $n$ 年后的量为 $b$，则：
\[
b = a(1+r)^n
\]

\textbf{各参数含义}：
\begin{itemize}
    \item $a$：基准量（增长前的原始数量）
    \item $r$：年平均增长率（以小数表示，如 $20\%$ 记为 $0.2$）
    \item $n$：增长年数（时间周期数）
    \item $b$：终值（增长后的最终数量）
\end{itemize}

\textbf{关键理解}："平均增长率"意味着每年的增长率相同，是复利增长模式。
\end{conceptbox}

\item \textbf{解题策略构建}

根据题意建立数学模型：
\begin{enumerate}[label=(\alph*)]
    \item \textbf{时间跨度}：从"去年"到"明年"共经过 \textbf{2年}
    \item \textbf{已知条件}：$a = 3$（去年），$b = 12$（明年），$n = 2$
    \item \textbf{所求量}：年平均增长率 $r$
    \item \textbf{方程建立}：代入公式 $12 = 3(1+r)^2$
\end{enumerate}

\item \textbf{详细求解过程}

\noindent\textbf{步骤1：建立方程}

根据增长率公式，设年平均增长率为 $x$（为避免与坐标变量混淆，这里用 $x$ 表示增长率）：
\[
3(1+x)^2 = 12
\]

\noindent\textbf{步骤2：化简方程}

两边同时除以 3：
\begin{align*}
(1+x)^2 &= \frac{12}{3} \\
(1+x)^2 &= 4
\end{align*}

\noindent\textbf{步骤3：开方求解}

两边同时开平方：
\begin{align*}
1+x &= \pm\sqrt{4} \\
1+x &= \pm 2
\end{align*}

得到两个方程：
\begin{itemize}
    \item 方程①：$1+x = 2$ \quad $\Rightarrow$ \quad $x = 1$
    \item 方程②：$1+x = -2$ \quad $\Rightarrow$ \quad $x = -3$
\end{itemize}

\noindent\textbf{步骤4：根据实际意义取舍}

\begin{itemize}
    \item $x = 1 = 100\%$：表示年平均增长率为 $100\%$，即每年资金翻一倍，\textbf{符合题意}
    \item $x = -3 = -300\%$：表示年平均减少率为 $300\%$，资金投入为负数，\textbf{不符合实际意义，舍去}
\end{itemize}

\item \textbf{验证与检验}

\noindent\textbf{验证计算}：
\begin{itemize}
    \item 去年投入：3 亿元
    \item 今年投入：$3 \times (1+100\%) = 3 \times 2 = 6$ 亿元
    \item 明年投入：$6 \times (1+100\%) = 6 \times 2 = 12$ 亿元 \checkmark
\end{itemize}

\noindent\textbf{逻辑检验}：
\begin{itemize}
    \item 增长率为正数，符合"逐年增加"的描述
    \item 两年正好从 3 亿增长到 12 亿（$3 \to 6 \to 12$），验证正确
\end{itemize}

\item \textbf{最终答案及表示方法}

\textbf{答案}：这两年中投入资金的年平均增长率为 \textbf{$\mathbf{100\%}$}（或写成 1）。

\textbf{完整答句}：

设这两年中投入资金的年平均增长率为 $x$，根据题意得：
\[
3(1+x)^2 = 12
\]
解得 $x_1 = 1 = 100\%$，$x_2 = -3$（不合题意，舍去）。

答：这两年中投入资金的年平均增长率为 $100\%$。
\end{detailedanalysis}

\subsection*{核心公式与推导}

\begin{formuladerivation}
\item \textbf{平均增长率公式}

对于平均增长率问题，基本公式为：
\[
a(1+r)^n = b
\]

\textbf{推导}：
\begin{itemize}
    \item 第1年末（经过1年）：$a(1+r)$
    \item 第2年末（经过2年）：$a(1+r)(1+r) = a(1+r)^2$
    \item 第 $n$ 年末（经过 $n$ 年）：$a(1+r)^n$
\end{itemize}

\item \textbf{平均下降率公式}

对于平均下降（减少）率问题，公式为：
\[
a(1-r)^n = b
\]

其中 $r$ 为平均下降率（$0 < r < 1$）。

\textbf{区别}：增长率用 $(1+r)$，下降率用 $(1-r)$。
\end{formuladerivation}

\subsection*{考点深度分析}

\begin{conceptbox}
\textbf{知识点归类}：
\begin{enumerate}
    \item \textbf{增长率概念}：平均增长率的数学含义
    \item \textbf{一元二次方程建模}：将实际问题转化为方程
    \item \textbf{解方程技巧}：直接开平方法
    \item \textbf{根的取舍}：根据实际意义判断解的合理性
    \item \textbf{复利思想}：指数增长模型
\end{enumerate}
\end{conceptbox}

\noindent\textbf{易错点全面解析}：
\begin{enumerate}[label=\arabic*.]
    \item \textbf{时间计算错误}：
    \begin{itemize}
        \item 错解：认为"去年到明年"是 1 年
        \item 正解：去年 $\to$ 今年 $\to$ 明年，共 \textbf{2年}
    \end{itemize}
    
    \item \textbf{方程列错}：
    \begin{itemize}
        \item 错解：$3(1+x) = 12$（一次方程）
        \item 正解：$3(1+x)^2 = 12$（二次方程）
    \end{itemize}
    
    \item \textbf{单位遗漏}：
    \begin{itemize}
        \item 错解：答"年平均增长率为 1"
        \item 正解：答"年平均增长率为 $100\%$"或"1（即 $100\%$）"
    \end{itemize}
    
    \item \textbf{未舍去负根}：
    \begin{itemize}
        \item 错解：保留 $x = -3$ 作为答案
        \item 正解：增长率为负数无实际意义，必须舍去
    \end{itemize}
    
    \item \textbf{开方漏解}：
    \begin{itemize}
        \item 错解：只取 $(1+x)^2 = 4$ 的正根 $1+x = 2$
        \item 正解：一元二次方程通常有两个根，都要考虑
    \end{itemize}
\end{enumerate}

\noindent\textbf{解题技巧与拓展}：
\begin{enumerate}[label=\arabic*.]
    \item \textbf{快速验证法}：
    
    对于增长率问题，验证时可用"倒推法"：
    \[
    \text{终值} \div (1+r) \div (1+r) = \text{初值}
    \]
    
    \item \textbf{常见题型归纳}：
    \begin{itemize}
        \item 类型1：资金/产值增长率（如本题）
        \item 类型2：人口增长率
        \item 类型3：产量/销量增长率
        \item 类型4：降价率问题（公式为 $a(1-r)^n = b$）
    \end{itemize}
    
    \item \textbf{拓展思考}：
    \begin{itemize}
        \item 如果题目改为"后年投入 24 亿元"，年平均增长率是多少？
        \item 如果年平均增长率为 $50\%$，3 年后投入是多少？
    \end{itemize}
\end{enumerate}
\end{redenv}
% ========== 第1题结束 ==========

\newpage

% ========== 第2题开始 ==========
\problemrule
\section*{【题目2】销售利润问题}
某店只销售某种进价为 40 元/kg 的特产. 已知该店按 60 元/kg 出售时，平均每天可售出 100 kg，后来经过市场调查发现，单价每降低 1 元，则平均每天的销售量可增加 10 kg. 若该店销售这种特产计划平均每天获利 2\,240 元.
\begin{enumerate}[label=(\arabic*)]
    \item 每千克该种特产应降价多少元？
    \item 为尽可能让利于顾客，则该店应按原售价的几折出售？
\end{enumerate}

\begin{redenv}
\analysisrule
\subsection*{逐步解析}

\begin{detailedanalysis}
\item \textbf{问题本质分析}

这是一道典型的\textbf{营销利润问题}，涉及多个变量的相互影响：
\begin{itemize}
    \item \textbf{进价}：固定成本，40 元/kg
    \item \textbf{售价}：变量，原售价 60 元/kg，可降价
    \item \textbf{销量}：变量，与售价相关，原价时 100 kg/天
    \item \textbf{利润}：目标值，2\,240 元/天
\end{itemize}

\begin{conceptbox}
\textbf{利润问题的核心概念}：

\textbf{基本公式}：
\[
\text{总利润} = \text{每千克利润} \times \text{销售数量}
\]

其中：
\[
\text{每千克利润} = \text{售价} - \text{进价}
\]

\textbf{变量关系}（本题关键）：
\begin{itemize}
    \item 单价每降低 1 元，销量增加 10 kg
    \item 这是一个\textbf{线性关系}：降价 $x$ 元 $\Rightarrow$ 销量增加 $10x$ kg
\end{itemize}
\end{conceptbox}

\item \textbf{解题策略构建}

\noindent\textbf{第(1)问策略}：

设每千克降价 $x$ 元，建立利润方程：
\begin{enumerate}[label=(\alph*)]
    \item 降价后售价：$(60 - x)$ 元/kg
    \item 降价后每千克利润：$(60 - x) - 40 = (20 - x)$ 元
    \item 降价后销量：$(100 + 10x)$ kg
    \item 总利润方程：$(20 - x)(100 + 10x) = 2240$
\end{enumerate}

\noindent\textbf{第(2)问策略}：

"尽可能让利于顾客"意味着要降价\textbf{最多}，即选择较大的降价金额，然后计算折扣。

\item \textbf{详细求解过程}

\noindent\textbf{第(1)问求解}：

设每千克该种特产应降价 $x$ 元。

根据题意，降价后：
\begin{itemize}
    \item 每千克售价：$(60 - x)$ 元
    \item 每千克利润：$(60 - x - 40) = (20 - x)$ 元
    \item 每天销售量：$(100 + 10x)$ kg
\end{itemize}

列方程：
\[
(20 - x)(100 + 10x) = 2240
\]

展开左边：
\begin{align*}
(20 - x)(100 + 10x) &= 20 \times 100 + 20 \times 10x - x \times 100 - x \times 10x \\
&= 2000 + 200x - 100x - 10x^2 \\
&= 2000 + 100x - 10x^2
\end{align*}

整理为标准形式：
\begin{align*}
2000 + 100x - 10x^2 &= 2240 \\
-10x^2 + 100x + 2000 - 2240 &= 0 \\
-10x^2 + 100x - 240 &= 0
\end{align*}

两边同除以 $-10$：
\[
x^2 - 10x + 24 = 0
\]

因式分解：
\[
(x - 4)(x - 6) = 0
\]

解得：
\[
x_1 = 4, \quad x_2 = 6
\]

\textbf{检验}：两个解都为正数，且降价后售价分别为 56 元和 54 元，都高于进价 40 元，\textbf{均符合题意}。

\noindent\textbf{第(2)问求解}：

"尽可能让利于顾客"意味着降价越多越好：
\begin{itemize}
    \item 选择 $x = 6$（降价 6 元）
    \item 降价后售价：$60 - 6 = 54$ 元
    \item 折扣率：$\dfrac{54}{60} = 0.9 = 90\%$
\end{itemize}

\item \textbf{验证与检验}

\noindent\textbf{验证 $x = 4$}：
\begin{itemize}
    \item 售价：$60 - 4 = 56$ 元/kg
    \item 每千克利润：$56 - 40 = 16$ 元
    \item 销量：$100 + 10 \times 4 = 140$ kg
    \item 总利润：$16 \times 140 = 2240$ 元 \checkmark
\end{itemize}

\noindent\textbf{验证 $x = 6$}：
\begin{itemize}
    \item 售价：$60 - 6 = 54$ 元/kg
    \item 每千克利润：$54 - 40 = 14$ 元
    \item 销量：$100 + 10 \times 6 = 160$ kg
    \item 总利润：$14 \times 160 = 2240$ 元 \checkmark
\end{itemize}

\noindent\textbf{验证折扣}：
\[
\frac{54}{60} \times 10 = 9 \text{（折）}
\]

\item \textbf{最终答案}

\textbf{答}：
\begin{enumerate}[label=(\arabic*)]
    \item 每千克该种特产应降价 \textbf{4 元或 6 元}。
    \item 为尽可能让利于顾客，该店应按原售价的 \textbf{九折}出售。
\end{enumerate}
\end{detailedanalysis}

\subsection*{核心公式与推导}

\begin{formuladerivation}
\item \textbf{利润基本公式}

\[
\text{总利润} = (\text{售价} - \text{进价}) \times \text{销售量}
\]

或记为：
\[
L = (p - c) \cdot q
\]

其中：$L$=利润，$p$=售价，$c$=进价，$q$=销售量。

\item \textbf{价格-销量关系式}

若单价每变化 1 元，销量变化 $k$ 单位，设降价 $x$ 元，则：
\[
\text{新售价} = p_0 - x
\]
\[
\text{新销量} = q_0 + kx
\]

其中 $p_0$、$q_0$ 为原售价和原销量。

\item \textbf{利润方程}

\[
(p_0 - x - c)(q_0 + kx) = L_{\text{目标}}
\]

展开后为一元二次方程。
\end{formuladerivation}

\subsection*{考点深度分析}

\begin{conceptbox}
\textbf{知识点归类}：
\begin{enumerate}
    \item \textbf{经济数学基础}：成本、售价、利润、销量的关系
    \item \textbf{函数思想}：利润关于降价金额的函数关系
    \item \textbf{一元二次方程建模}：根据等量关系列方程
    \item \textbf{因式分解法}：解一元二次方程
    \item \textbf{实际意义判断}：双根情况下如何选择
\end{enumerate}
\end{conceptbox}

\noindent\textbf{易错点全面解析}：
\begin{enumerate}[label=\arabic*.]
    \item \textbf{利润计算错误}：
    \begin{itemize}
        \item 错解：用售价代替利润列方程 $60 \times (100 + 10x) = 2240$
        \item 正解：利润 = 售价 - 进价，应为 $(60 - x - 40)(100 + 10x) = 2240$
    \end{itemize}
    
    \item \textbf{销量变化方向错误}：
    \begin{itemize}
        \item 错解：降价 $x$ 元，销量减少 $10x$
        \item 正解：降价促进销售，销量\textbf{增加}$10x$
    \end{itemize}
    
    \item \textbf{单位不统一}：
    \begin{itemize}
        \item 注意：进价、售价、降价金额单位要一致（都是元）
    \end{itemize}
    
    \item \textbf{第(2)问理解错误}：
    \begin{itemize}
        \item 错解：选择降价 4 元（较小的降价）
        \item 正解："让利于顾客"意味着降价更多，应选 6 元
    \end{itemize}
    
    \item \textbf{折扣计算错误}：
    \begin{itemize}
        \item 错解：$\frac{6}{60} = 10\%$（打一折）
        \item 正解：折扣 $= \frac{\text{现价}}{\text{原价}} = \frac{54}{60} = 0.9$（九折）
    \end{itemize}
\end{enumerate}
\end{redenv}
% ========== 第2题结束 ==========

\newpage

% ========== 第3题开始 ==========
\problemrule
\section*{【题目3】面积问题（平移法）}
某单位准备将院内一块长 30 m、宽 20 m 的长方形空地，建成一个矩形花园，要求在花园中修建两条纵向和一条横向的小道，剩余的地方种植花草，如图所示. 要使种植花草的面积为 532 m$^2$，那么小道进出口的宽度应为多少米？（注：所有小道进出口的宽度相等.）

\begin{center}
\begin{tikzpicture}[scale=0.12]
    \def\W{30}
    \def\H{20}
    \def\pw{2}
    \fill[grassgreen] (0,0) rectangle (\W, \H);
    \filldraw[fill=gray!20, draw=black] (0, \H/2 - \pw/2) rectangle (\W, \H/2 + \pw/2);
    \filldraw[fill=gray!20, draw=black] (\W/3 - \pw/2, 0) rectangle (\W/3 + \pw/2, \H);
    \filldraw[fill=gray!20, draw=black] (2*\W/3 - \pw/2, 0) rectangle (2*\W/3 + \pw/2, \H);
    \draw[thick] (0,0) rectangle (\W, \H);
    \draw[<->, >=latex] (0, -2) -- (\W, -2) node[midway, fill=white] {30 m};
    \draw[<->, >=latex] (\W+2, 0) -- (\W+2, \H) node[midway, fill=white] {20 m};
    \node at (\W/2, -5) {(第 3 题图)};
\end{tikzpicture}
\end{center}

\begin{redenv}
\analysisrule
\subsection*{逐步解析}

\begin{detailedanalysis}
\item \textbf{问题本质分析}

这是一道\textbf{几何面积问题}，涉及在矩形区域内修建小道后剩余种植面积的计算。

\begin{conceptbox}
\textbf{核心思路——平移法}：

小道将花园分割成若干块，直接计算种植面积较复杂。

\textbf{平移法原理}：将分散的种植区域通过平移拼接成一个新的矩形。
\begin{itemize}
    \item 纵向小道宽度为 $x$，则种植区域的有效宽度为 $(30 - 2x)$ m
    \item 横向小道宽度为 $x$，则种植区域的有效高度为 $(20 - x)$ m
    \item 种植总面积 = $(30 - 2x)(20 - x)$
\end{itemize}

\textbf{关键理解}：平移不改变面积，但可以将复杂的分散图形转化为简单矩形。
\end{conceptbox}

\item \textbf{解题策略构建}

设小道进出口的宽度为 $x$ 米。

\begin{enumerate}[label=(\alph*)]
    \item \textbf{平移后的有效尺寸}：
    \begin{itemize}
        \item 原宽度：30 m，扣除两条纵向小道：$30 - 2x$
        \item 原高度：20 m，扣除一条横向小道：$20 - x$
    \end{itemize}
    
    \item \textbf{种植面积方程}：
    \[
    (30 - 2x)(20 - x) = 532
    \]
    
    \item \textbf{约束条件}：
    \begin{itemize}
        \item $x > 0$（小道必须有宽度）
        \item $30 - 2x > 0$ 即 $x < 15$（纵向不能超出）
        \item $20 - x > 0$ 即 $x < 20$（横向不能超出）
    \end{itemize}
\end{enumerate}

\item \textbf{详细求解过程}

设小道进出口的宽度为 $x$ 米。

根据平移法，种植区域的有效尺寸为：
\begin{itemize}
    \item 宽度：$(30 - 2x)$ 米
    \item 高度：$(20 - x)$ 米
\end{itemize}

列方程：
\[
(30 - 2x)(20 - x) = 532
\]

展开左边：
\begin{align*}
(30 - 2x)(20 - x) &= 30 \times 20 - 30x - 40x + 2x^2 \\
&= 600 - 70x + 2x^2
\end{align*}

整理为标准形式：
\begin{align*}
600 - 70x + 2x^2 &= 532 \\
2x^2 - 70x + 600 - 532 &= 0 \\
2x^2 - 70x + 68 &= 0
\end{align*}

两边同除以 2：
\[
x^2 - 35x + 34 = 0
\]

因式分解：
\[
(x - 1)(x - 34) = 0
\]

解得：
\[
x_1 = 1, \quad x_2 = 34
\]

\textbf{检验}：
\begin{itemize}
    \item $x = 34$：$30 - 2 \times 34 = 30 - 68 = -38 < 0$，不符合实际意义，舍去
    \item $x = 1$：$30 - 2 \times 1 = 28 > 0$，$20 - 1 = 19 > 0$，符合题意
\end{itemize}

\item \textbf{验证与检验}

当 $x = 1$ 时：
\begin{itemize}
    \item 有效宽度：$30 - 2 \times 1 = 28$ m
    \item 有效高度：$20 - 1 = 19$ m
    \item 种植面积：$28 \times 19 = 532$ m$^2$ \checkmark
\end{itemize}

\item \textbf{最终答案}

\textbf{答}：小道进出口的宽度应为 \textbf{1 米}。
\end{detailedanalysis}

\subsection*{核心公式与推导}

\begin{formuladerivation}
\item \textbf{平移法面积公式}

对于在 $a \times b$ 的矩形中修建宽度为 $x$ 的纵横小道：

若有 $m$ 条纵向小道和 $n$ 条横向小道，则：
\[
\text{种植面积} = (a - mx)(b - nx)
\]

\textbf{推导}：将种植区域平移拼接后，形成新矩形的两边分别为原长减去小道总宽度。

\item \textbf{面积方程标准形式}

\[
(a - mx)(b - nx) = S
\]

展开后：
\[
mnx^2 - (an + bm)x + (ab - S) = 0
\]

这是一元二次方程的标准形式。
\end{formuladerivation}

\subsection*{考点深度分析}

\begin{conceptbox}
\textbf{知识点归类}：
\begin{enumerate}
    \item \textbf{几何直观}：理解小道的分布和种植区域的形状
    \item \textbf{平移思想}：将复杂图形转化为简单矩形
    \item \textbf{多项式乘法}：展开方程
    \item \textbf{一元二次方程求解}：因式分解法
    \item \textbf{实际约束}：根据几何意义筛选解
\end{enumerate}
\end{conceptbox}

\noindent\textbf{易错点全面解析}：
\begin{enumerate}[label=\arabic*.]
    \item \textbf{重复扣除}：
    \begin{itemize}
        \item 错解：直接计算 $30 \times 20 - $ 小道面积（小道交叉处被重复扣除）
        \item 正解：使用平移法避免重复计算
    \end{itemize}
    
    \item \textbf{小道数量错误}：
    \begin{itemize}
        \item 错解：$30 - x$（只减一条纵向小道）
        \item 正解：$30 - 2x$（两条纵向小道）
    \end{itemize}
    
    \item \textbf{未舍去不合理解}：
    \begin{itemize}
        \item $x = 34$ 会使有效宽度为负数，必须舍去
    \end{itemize}
    
    \item \textbf{方程展开错误}：
    \begin{itemize}
        \item 注意符号：$(-2x) \times (-x) = +2x^2$
    \end{itemize}
\end{enumerate}
\end{redenv}
% ========== 第3题结束 ==========

\newpage

% ========== 第4题开始 ==========
\problemrule
\section*{【题目4】日历表问题（数列规律）}
如图是某年 8 月的日历表，在此日历表上可以用一个矩形圈出 $3 \times 3$ 个位置相邻的 9 个数（如 6, 7, 8, 13, 14, 15, 20, 21, 22）。若圈出的 9 个数的和为 117，求这 9 个数中最小的数。

\begin{center}
\begin{minipage}{0.55\textwidth}
\centering
\renewcommand{\arraystretch}{1.3}
\setlength{\tabcolsep}{5pt}
\begin{tabular}{|c|c|c|c|c|c|c|}
    \hline
    \rowcolor{black!10}
    \textcolor{weekendred}{日} & 一 & 二 & 三 & 四 & 五 & \textcolor{weekendred}{六} \\
    \hline
     & & & 1 & 2 & 3 & \textcolor{weekendred}{4} \\
     & & & \scriptsize 建军节 & \scriptsize 十五 & \scriptsize 十六 & \textcolor{weekendred}{\scriptsize 十七} \\
    \hline
    \textcolor{weekendred}{5} & 6 & 7 & 8 & 9 & 10 & \textcolor{weekendred}{11} \\
    \textcolor{weekendred}{\scriptsize 十八} & \scriptsize 十九 & \textcolor{solargreen}{\scriptsize 立秋} & \scriptsize 廿一 & \scriptsize 廿二 & \scriptsize 廿三 & \textcolor{weekendred}{\scriptsize 廿四} \\
    \hline
    \textcolor{weekendred}{12} & 13 & 14 & 15 & 16 & 17 & \textcolor{weekendred}{18} \\
    \textcolor{weekendred}{\scriptsize 廿五} & \scriptsize 廿六 & \scriptsize 廿七 & \scriptsize 廿八 & \scriptsize 廿九 & \scriptsize 七月 & \textcolor{weekendred}{\scriptsize 初二} \\
    \hline
    \textcolor{weekendred}{19} & 20 & 21 & 22 & 23 & 24 & \textcolor{weekendred}{25} \\
    \textcolor{weekendred}{\scriptsize 初三} & \scriptsize 初四 & \scriptsize 初五 & \scriptsize 初六 & \scriptsize 初七 & \scriptsize 初八 & \textcolor{weekendred}{\scriptsize 初九} \\
    \hline
    \textcolor{weekendred}{26} & 27 & 28 & 29 & 30 & 31 & \\
    \textcolor{weekendred}{\scriptsize 初十} & \scriptsize 十一 & \scriptsize 十二 & \scriptsize 十三 & \scriptsize 十四 & \scriptsize 十五 & \\
    \hline
\end{tabular}
\smallskip
\small{(第 4 题图)}
\end{minipage}
\end{center}

\begin{redenv}
\analysisrule
\subsection*{逐步解析}

\begin{detailedanalysis}
\item \textbf{问题本质分析}

这是一道\textbf{数字规律问题}，需要发现日历表中 $3 \times 3$ 矩形框内 9 个数的排列规律。

\begin{conceptbox}
\textbf{日历表的数字规律}：

在日历表中，相邻数字之间存在固定关系：
\begin{itemize}
    \item \textbf{横向}：左右相邻的数相差 1
    \item \textbf{纵向}：上下相邻的数相差 7（一周有 7 天）
\end{itemize}

对于 $3 \times 3$ 矩形框中的 9 个数，设最小数（左上角）为 $x$，则：

\[
\begin{array}{ccc}
x & x+1 & x+2 \\
x+7 & x+8 & x+9 \\
x+14 & x+15 & x+16
\end{array}
\]
\end{conceptbox}

\item \textbf{解题策略构建}

设 $3 \times 3$ 矩形框中最小的数为 $x$，则 9 个数分别为：
\begin{itemize}
    \item 第一行：$x$, $x+1$, $x+2$
    \item 第二行：$x+7$, $x+8$, $x+9$
    \item 第三行：$x+14$, $x+15$, $x+16$
\end{itemize}

根据题意，这 9 个数的和为 117，列方程求解。

\item \textbf{详细求解过程}

设这 9 个数中最小的数为 $x$。

根据日历规律，这 9 个数为：
\[
\begin{array}{ccc}
x & x+1 & x+2 \\
x+7 & x+8 & x+9 \\
x+14 & x+15 & x+16
\end{array}
\]

9 个数之和：
\begin{align*}
S &= x + (x+1) + (x+2) + (x+7) + (x+8) + (x+9) + (x+14) + (x+15) + (x+16) \\
&= 9x + (1+2+7+8+9+14+15+16) \\
&= 9x + 72
\end{align*}

根据题意：
\begin{align*}
9x + 72 &= 117 \\
9x &= 45 \\
x &= 5
\end{align*}

\textbf{更快捷的方法}——利用中心数性质：

在 $3 \times 3$ 矩形框中，\textbf{9 个数的和 = 中心数 $\times$ 9}。

\begin{itemize}
    \item 9 个数之和为 117
    \item 中心数为 $117 \div 9 = 13$
    \item 最小数比中心数小 8，所以最小数为 $13 - 8 = 5$
\end{itemize}

\item \textbf{验证与检验}

当最小数为 5 时，9 个数为：
\[
\begin{array}{ccc}
5 & 6 & 7 \\
12 & 13 & 14 \\
19 & 20 & 21
\end{array}
\]

求和：
\[
5+6+7+12+13+14+19+20+21 = 117 \quad \checkmark
\]

从日历表上看，这 9 个数正好位于 8 月 5 日（周日）到 8 月 21 日（周二），在日历表中形成一个完整的 $3 \times 3$ 矩形框。

\item \textbf{最终答案}

\textbf{答}：这 9 个数中最小的数是 \textbf{5}。
\end{detailedanalysis}

\subsection*{核心公式与推导}

\begin{formuladerivation}
\item \textbf{$3 \times 3$ 矩形框求和公式}

对于日历表中以 $x$ 为最小数的 $3 \times 3$ 矩形框：
\[
S = 9x + 72
\]
或等价地：
\[
S = 9(x + 8) = 9 \times \text{中心数}
\]

\textbf{推导}：9 个数构成等差数列的变形，中心数为 $x+8$，9 个数以中心数对称分布。

\item \textbf{通用规律}

对于任意 $n \times n$ 的日历矩形框，数字之和与中心数（或中心位置）存在确定关系。
\end{formuladerivation}

\subsection*{考点深度分析}

\begin{conceptbox}
\textbf{知识点归类}：
\begin{enumerate}
    \item \textbf{数列规律}：发现日历中数字的排列规律
    \item \textbf{代数建模}：用字母表示最小数，建立方程
    \item \textbf{简化技巧}：利用中心数性质快速求解
    \item \textbf{整除性质}：9 个数的和必能被 9 整除
\end{enumerate}
\end{conceptbox}

\noindent\textbf{解题技巧}：
\begin{enumerate}
    \item \textbf{中心数法}：$3 \times 3$ 框中，和 $= 9 \times$ 中心数（最快捷）
    \item \textbf{平均数法}：9 个数的平均数就是中心数
\end{enumerate}
\end{redenv}
% ========== 第4题结束 ==========

\newpage

% ========== 第7题开始 ==========
\problemrule
\section*{【题目7】动点问题（平行四边形判定）}
如图，在矩形 $ABCD$ 中，$BC = 24$ cm，$P, Q, M, N$ 分别从点 $A, B, C, D$ 同时出发，分别沿边 $AD, BC, CB, DA$ 移动，且当有一个点先到达所在边的另一个端点时，其他各点也随之停止移动. 已知移动一段时间后，若 $BQ = x$ cm（$x \neq 0$），则 $AP = 2x$ cm，$CM = 3x$ cm，$DN = x^2$ cm.

\begin{center}
\begin{tikzpicture}[scale=0.35]
    \coordinate (A) at (0, 6);
    \coordinate (B) at (0, 0);
    \coordinate (C) at (12, 0);
    \coordinate (D) at (12, 6);
    \draw[thick] (A) -- (B) -- (C) -- (D) -- cycle;
    \node[left] at (A) {$A$};
    \node[left] at (B) {$B$};
    \node[right] at (C) {$C$};
    \node[right] at (D) {$D$};
    \coordinate (P) at (2, 6);
    \coordinate (N) at (10, 6);
    \coordinate (Q) at (1.5, 0);
    \coordinate (M) at (9, 0);
    \fill (P) circle (3pt) node[above] {$P$};
    \fill (N) circle (3pt) node[above] {$N$};
    \fill (Q) circle (3pt) node[below] {$Q$};
    \fill (M) circle (3pt) node[below] {$M$};
    \draw[thick] (P) -- (Q);
    \draw[thick] (N) -- (M);
    \draw[->, >=stealth, bookblue, thick] (A) -- (1, 6);
    \draw[->, >=stealth, bookblue, thick] (D) -- (11, 6);
    \draw[->, >=stealth, bookblue, thick] (B) -- (0.75, 0);
    \draw[->, >=stealth, bookblue, thick] (C) -- (10.5, 0);
    \node at (6, -2) {(第 7 题图)};
\end{tikzpicture}
\end{center}

\begin{enumerate}[label=(\arabic*)]
    \item 当 $x$ 为何值时，$P, N$ 两点重合？
    \item 问 $Q, M$ 两点能重合吗？若 $Q, M$ 两点能重合，则求出相应的 $x$ 的值；若 $Q, M$ 两点不能重合，请说明理由.
    \item 当 $x$ 为何值时，以 $P, Q, M, N$ 为顶点的四边形是平行四边形？
\end{enumerate}

\begin{redenv}
\analysisrule
\subsection*{逐步解析}

\begin{detailedanalysis}
\item \textbf{问题本质分析}

这是一道\textbf{动态几何问题}，涉及矩形中的动点运动和平行四边形判定。

\begin{conceptbox}
\textbf{关键几何关系}：

首先确定矩形各边长度：
\begin{itemize}
    \item $BC = 24$ cm（已知）
    \item 在矩形中，$AD = BC = 24$ cm，$AB = CD$（具体值未知，但不影响本题）
\end{itemize}

\textbf{各点位置表示}（沿边移动的距离）：
\begin{itemize}
    \item $P$ 从 $A$ 出发沿 $AD$ 移动：$AP = 2x$，到 $D$ 的距离 $PD = 24 - 2x$
    \item $Q$ 从 $B$ 出发沿 $BC$ 移动：$BQ = x$
    \item $M$ 从 $C$ 出发沿 $CB$ 移动：$CM = 3x$，到 $B$ 的距离 $MB = 24 - 3x$
    \item $N$ 从 $D$ 出发沿 $DA$ 移动：$DN = x^2$，到 $A$ 的距离 $NA = 24 - x^2$
\end{itemize}

\textbf{停止条件}：任意一点到达端点时，所有点停止。
\end{conceptbox}

\item \textbf{解题策略构建}

\noindent\textbf{第(1)问——$P, N$ 重合}：

$P$ 和 $N$ 都在边 $AD$ 上，两点重合意味着它们到 $A$ 点的距离相等：
\[
AP = AN \quad \text{即} \quad AP = AD - DN
\]

\noindent\textbf{第(2)问——$Q, M$ 重合}：

$Q$ 和 $M$ 都在边 $BC$ 上（$Q$ 从 $B$ 向 $C$，$M$ 从 $C$ 向 $B$），重合条件：
\[
BQ + CM = BC \quad \text{或} \quad BQ = BM \quad \text{（即 $x = 24 - 3x$）}
\]

\noindent\textbf{第(3)问——平行四边形判定}：

连接 $PQ$、$QM$、$MN$、$NP$ 形成四边形。

由图可知，$PQ$ 和 $NM$ 是两条连线。若 $PQ \parallel NM$ 且 $PQ = NM$，则四边形为平行四边形。

由于 $P, N$ 在 $AD$ 上，$Q, M$ 在 $BC$ 上，且 $AD \parallel BC$，所以只要 $PN$ 和 $QM$ 的中点重合（或对角线互相平分），即可判定为平行四边形。

更简便的方法：由于 $PQNM$ 的顶点分别在平行的两边上，只需满足 $PN = QM$（上下两边平行且相等）。

注意：$PN = |AP - AN|$ 或 $PN = |AD - AP - DN|$，具体取决于位置。

实际上，$P$ 在 $A$ 和 $N$ 之间还是 $N$ 在 $A$ 和 $P$ 之间需要讨论。

\item \textbf{详细求解过程}

\noindent\textbf{第(1)问求解}：

$P, N$ 两点重合 $\Leftrightarrow AP = AN$

即：
\begin{align*}
AP &= AD - DN \\
2x &= 24 - x^2 \\
x^2 + 2x - 24 &= 0
\end{align*}

因式分解：
\[
(x + 6)(x - 4) = 0
\]

解得：$x = -6$（舍去，因为 $x > 0$）或 $x = 4$

\textbf{检验}：当 $x = 4$ 时，$AP = 8$，$DN = 16$，$AN = 24 - 16 = 8$，符合。

\noindent\textbf{第(2)问求解}：

$Q, M$ 两点重合 $\Leftrightarrow BQ + CM = BC$（因为它们从两端相向而行）

即：
\begin{align*}
x + 3x &= 24 \\
4x &= 24 \\
x &= 6
\end{align*}

\textbf{检验停止条件}：
\begin{itemize}
    \item $AP = 2 \times 6 = 12 < 24$（未到 $D$）
    \item $DN = 36 > 24$（已经超过 $A$！）
\end{itemize}

由于当 $x = 6$ 时，$DN = 36 > 24$，说明点 $N$ 早已到达 $A$ 点并停止，所以实际上 $x$ 不能取到 6。

更准确地说，$N$ 点到达 $A$ 点时：$DN = AD = 24$，即 $x^2 = 24$，$x = \sqrt{24} = 2\sqrt{6} \approx 4.9$

因此当 $x = 6$ 时，运动早已停止，$Q, M$ 两点\textbf{不能重合}。

\noindent\textbf{第(3)问求解}：

四边形 $PQNM$ 为平行四边形的条件是 $PN = QM$（且方向相同）。

首先确定各点在边 $AD$ 和 $BC$ 上的位置：

\textbf{在边 $AD$ 上}：
\begin{itemize}
    \item $P$ 距 $A$ 点：$AP = 2x$
    \item $N$ 距 $D$ 点：$DN = x^2$，距 $A$ 点：$AN = 24 - x^2$
\end{itemize}

$PN$ 的长度需要分情况：
\begin{itemize}
    \item 若 $P$ 在 $N$ 左侧（靠近 $A$）：$PN = AN - AP = 24 - x^2 - 2x$（要求 $AN > AP$，即 $24 - x^2 > 2x$）
    \item 若 $N$ 在 $P$ 左侧：$PN = AP - AN = 2x - (24 - x^2) = x^2 + 2x - 24$
\end{itemize}

\textbf{在边 $BC$ 上}：
\begin{itemize}
    \item $Q$ 距 $B$ 点：$BQ = x$
    \item $M$ 距 $B$ 点：$BM = 24 - 3x$
\end{itemize}

$QM$ 的长度：$QM = |BM - BQ| = |24 - 3x - x| = |24 - 4x|$

对于平行四边形，需要 $PN = QM$。

\textbf{情况1}：$P$ 在 $N$ 左侧（$24 - x^2 > 2x$，即 $x^2 + 2x - 24 < 0$，$0 < x < 4$）

此时 $PN = 24 - x^2 - 2x$

\begin{itemize}
    \item 若 $24 - 4x \geq 0$（即 $x \leq 6$）：$QM = 24 - 4x$
    
    方程：$24 - x^2 - 2x = 24 - 4x$
    \begin{align*}
    -x^2 - 2x &= -4x \\
    -x^2 + 2x &= 0 \\
    x(-x + 2) &= 0
    \end{align*}
    
    $x = 0$（舍去）或 $x = 2$
    
    \item 验证：$x = 2$ 在 $(0, 4)$ 范围内，有效。
\end{itemize}

\textbf{情况2}：$N$ 在 $P$ 左侧（$x > 4$）

此时 $PN = x^2 + 2x - 24$

\begin{itemize}
    \item 若 $x < 6$：$QM = 24 - 4x$
    
    方程：$x^2 + 2x - 24 = 24 - 4x$
    \begin{align*}
    x^2 + 6x - 48 &= 0
    \end{align*}
    
    求根公式：$x = \dfrac{-6 + \sqrt{36 + 192}}{2} = \dfrac{-6 + \sqrt{228}}{2} = -3 + \sqrt{57} \approx 4.55$
    
    在 $(4, 6)$ 范围内，有效。
    
    \item 若 $x > 6$：$QM = 4x - 24$
    
    方程：$x^2 + 2x - 24 = 4x - 24$
    \begin{align*}
    x^2 - 2x &= 0 \\
    x(x - 2) &= 0
    \end{align*}
    
    $x = 0$（舍去）或 $x = 2$（不在 $x > 6$ 范围，舍去）
\end{itemize}

\textbf{约束条件检查}：

所有点都不能超出边界，即：
\begin{itemize}
    \item $AP = 2x \leq 24$，即 $x \leq 12$
    \item $BQ = x \leq 24$
    \item $CM = 3x \leq 24$，即 $x \leq 8$
    \item $DN = x^2 \leq 24$，即 $x \leq \sqrt{24} \approx 4.9$
\end{itemize}

最严格的约束是 $x \leq \sqrt{24} \approx 4.9$。

因此只有 $x = 2$ 和 $x = -3 + \sqrt{57} \approx 4.55$ 满足条件。

\item \textbf{最终答案}

\textbf{答}：
\begin{enumerate}[label=(\arabic*)]
    \item 当 $x = $\textbf{4} 时，$P, N$ 两点重合。
    
    \item $Q, M$ 两点\textbf{不能重合}。因为若 $Q, M$ 重合，则 $x = 6$，此时 $DN = 36 > 24$，点 $N$ 早已到达端点停止移动，矛盾。
    
    \item 当 $x = $\textbf{2} 或 $x = \sqrt{57} - 3$（约 4.55）时，以 $P, Q, M, N$ 为顶点的四边形是平行四边形。
\end{enumerate}
\end{detailedanalysis}

\subsection*{核心公式与推导}

\begin{formuladerivation}
\item \textbf{平行四边形判定定理}

在平面几何中，四边形为平行四边形的判定条件：
\begin{enumerate}
    \item 两组对边分别平行
    \item 两组对边分别相等
    \item 对角线互相平分
    \item 一组对边平行且相等
\end{enumerate}

\item \textbf{动点距离公式}

对于沿同一直线相向运动的两个点：
\[
\text{两点间距离} = |\text{路程}_1 - \text{路程}_2| \quad \text{或} \quad \text{总路程} - (\text{路程}_1 + \text{路程}_2)
\]
（取决于相对位置）
\end{formuladerivation}

\subsection*{考点深度分析}

\begin{conceptbox}
\textbf{知识点归类}：
\begin{enumerate}
    \item \textbf{矩形性质}：对边相等且平行
    \item \textbf{动点问题}：用代数式表示动点位置
    \item \textbf{平行四边形判定}：几何图形的判定条件
    \item \textbf{分类讨论}：不同情况下的位置关系
    \item \textbf{约束条件}：实际问题中的取值范围限制
\end{enumerate}
\end{conceptbox}

\noindent\textbf{易错点全面解析}：
\begin{enumerate}[label=\arabic*.]
    \item \textbf{忽略停止条件}：
    \begin{itemize}
        \item 必须检查 $DN = x^2 \leq 24$，这是最容易超出的限制
    \end{itemize}
    
    \item \textbf{$PN$ 长度计算错误}：
    \begin{itemize}
        \item 需要分 $P$ 在 $N$ 左侧还是右侧讨论
    \end{itemize}
    
    \item \textbf{$QM$ 长度符号错误}：
    \begin{itemize}
        \item 注意用绝对值表示，或分情况讨论
    \end{itemize}
    
    \item \textbf{遗漏解}：
    \begin{itemize}
        \item 平行四边形问题通常有多个解，要全面考虑
    \end{itemize}
\end{enumerate}
\end{redenv}
% ========== 第7题结束 ==========

\vspace{2em}
\begin{center}
\rule{0.6\textwidth}{0.5pt}\\[0.5em]
{\large\bfseries 讲义结束}
\end{center}

\end{document}
